\chapter{Ehlers-Danlos Syndrome}

\section*{Definition and Overview}
Ehlers-Danlos syndrome (EDS) is a group of inherited conditions that affect the connective tissue, which provides structure and support to the skin, joints, blood vessels, and other organs. The hallmark features are overly flexible joints, stretchy skin, and fragile tissues that bruise and scar easily. There are several types of EDS, ranging from the relatively common hypermobility type to rare, serious vascular forms. There is no cure, but most types are managed well with physiotherapy, pain management, and careful attention to skin and joint health.

\section*{Epidemiology}
The most common type, hypermobile EDS, affects an estimated 1 in 5,000 to 20,000 people, while the other types are much rarer. EDS affects people of all ethnicities and both sexes, and symptoms often become apparent in childhood with flexible joints and frequent sprains. Because mild joint hypermobility is common in the general population, EDS can be under-recognised, and many people wait years for a diagnosis.

\section*{Aetiology and Pathophysiology}
EDS is caused by changes in the genes that produce collagen and related proteins, the building blocks of connective tissue.
\begin{itemize}
    \item \textbf{Collagen defects:} mutations in collagen genes or the proteins that process collagen weaken the connective tissue
    \item \textbf{Types:} the major types include hypermobile EDS, classical EDS, and vascular EDS, each with different genes and features
    \item \textbf{Inheritance:} most types are inherited in an autosomal dominant pattern, so a child of an affected parent has a 50 per cent chance of inheriting it; new mutations also occur
    \item \textbf{Effects on the body:} weak connective tissue affects joints, causing laxity and dislocations; skin, causing stretchiness and fragility; and, in vascular EDS, blood vessels and organs
    \item \textbf{Hypermobility:} joints move beyond the normal range, leading to instability, pain, and injury
\end{itemize}

\section*{Clinical Features}
The features vary by type and severity.
\begin{itemize}
    \item Joint hypermobility, with frequent dislocations and sprains
    \item Chronic joint and muscle pain
    \item Stretchy, soft, velvety skin
    \item Easy bruising and slow-healing, widened scars
    \item Fatigue and muscle weakness
    \item Flat feet, scoliosis, and other musculoskeletal issues
    \item Digestive symptoms, including reflux and irritable bowel-like symptoms
    \item Dizziness and palpitations on standing, related to autonomic dysfunction in some people
    \item In vascular EDS: fragile blood vessels and organs, with risk of rupture
\end{itemize}

\section*{Differential Diagnosis}
Other conditions cause joint laxity and similar symptoms.
\begin{itemize}
    \item Benign joint hypermobility syndrome, now considered part of the same spectrum
    \item Marfan syndrome and other connective tissue disorders
    \item Hypermobility in young people without connective tissue disease
    \item Chronic pain conditions, including fibromyalgia
    \item Other genetic syndromes affecting the joints and skin
\end{itemize}

\section*{When to Seek Medical Care}
People with very flexible joints, frequent dislocations, stretchy or fragile skin, or unexplained bruising should be assessed by a GP and referred to a rheumatologist, geneticist, or specialist service. Families with a history of EDS should seek genetic counselling.

\textbf{Call 000 (or local emergency services) for:}
\begin{itemize}
    \item Sudden severe abdominal or chest pain, which in vascular EDS may indicate rupture of an organ or blood vessel
    \item Sudden severe headache or stroke symptoms
    \item Fainting or collapse
    \item A dislocated joint that will not move or is numb
\end{itemize}

\section*{Diagnosis and Investigations}
Diagnosis is made clinically and confirmed with genetic testing where appropriate.
\begin{itemize}
    \item \textbf{Clinical assessment:} evaluation of joint mobility, skin features, and family history against the diagnostic criteria
    \item \textbf{Beighton score:} a standard measure of joint hypermobility
    \item \textbf{Genetic testing:} blood tests identify the gene changes in the classical, vascular, and other types
    \item \textbf{Cardiac assessment:} echocardiography may be recommended because some types affect the heart and aorta
    \item \textbf{Specialist review:} rheumatology, genetics, and physiotherapy assessments
\end{itemize}

\section*{Management}
Management focuses on symptom control and preventing injury.
\begin{itemize}
    \item \textbf{Physiotherapy:} the cornerstone of care, building muscle strength to stabilise joints
    \item \textbf{Pain management:} a multidisciplinary approach, including exercise, pacing, and medicines where needed
    \item \textbf{Joint care:} bracing, supports, and techniques to prevent dislocations
    \item \textbf{Skin care:} protecting the skin, using sunscreen, and promoting careful wound healing
    \item \textbf{Activity planning:} low-impact exercise such as swimming and cycling
    \item \textbf{Managing autonomic symptoms:} fluids, salt, and strategies for dizziness
    \item \textbf{Specific care for vascular EDS:} specialist monitoring, blood pressure control, and avoidance of high-risk activities
    \item \textbf{Surgery:} considered carefully because tissues may heal poorly
\end{itemize}

\section*{Complications}
\begin{itemize}
    \item Recurrent joint dislocations and chronic pain
    \item Osteoarthritis from joint instability
    \item Poor wound healing and widened scars
    \item Chronic fatigue and reduced activity
    \item In vascular EDS: life-threatening rupture of blood vessels or organs
    \item Pregnancy complications in some types, requiring specialist care
\end{itemize}

\section*{Prognosis and Outcomes}
The outlook depends on the type of EDS. Hypermobile and classical types are not life-threatening, and with good physiotherapy and pain management most people maintain good function and a normal lifespan. Vascular EDS carries serious risks of vessel and organ rupture, but specialist care and monitoring improve outcomes. Genetic diagnosis allows families to understand the condition and plan appropriately.

\section*{Prevention and Self-Care}
\begin{itemize}
    \item Strengthen muscles around unstable joints with guided physiotherapy
    \item Choose low-impact activities and pace yourself
    \item Protect skin from sun and injury
    \item Seek early treatment of joint injuries
    \item Manage stress and fatigue with rest and good sleep
    \item Join support organisations for information and community
    \item Attend regular specialist reviews
\end{itemize}

\section*{Sources and Further Reading}
The following reputable sources provide further information for healthcare students and patients seeking reliable, current advice about Ehlers-Danlos syndrome.
\begin{itemize}
    \item Ehlers-Danlos Syndrome Australasia. \textit{Information and support for people with EDS.} https://edsa.org.au/
    \item The Ehlers-Danlos Society. \textit{Diagnosis and resources.} https://www.ehlers-danlos.com/
    \item NHS. \textit{Ehlers-Danlos syndromes: overview.} https://www.nhs.uk/conditions/ehlers-danlos-syndromes/
    \item Mayo Clinic. \textit{Ehlers-Danlos syndrome: symptoms and causes.} https://www.mayoclinic.org/diseases-conditions/ehlers-danlos-syndrome/
    \item MedlinePlus. \textit{Ehlers-Danlos syndrome.} https://medlineplus.gov/ehlersdanlossyndrome.html
    \item Genetic Alliance Australia. \textit{Genetic condition information.} https://www.geneticalliance.org.au/
\end{itemize}
